\subsection{Odd octic characters and initialization on the cubic source}

\begin{statement}{A.125}{Proved}
{All odd octic evaluations and finite cubic-source initialization}
Let $c_1,\ldots,c_k\in\Q$ be distinct, let
$a_i\in\Q^\times$, and consider the actual old cover
\[
              C:\quad T-c_i=a_iX_i^8\quad(1\le i\le k).
\tag{A.125.1}
\]
For $c\in\Q\setminus\{c_i\}$, let $B_8$ be the common
abelian field of the boundary fibre with radicands
$D_i=(c-c_i)/a_i$, computed in A.121.
For every character $\chi$ of $\operatorname{Gal}(B_8/\Q)$,
the evaluation of $(T-c,\chi)$ is constant on the ordinary
points of $C(\Q_p)$ avoiding $T=c$, at every odd prime $p$.
The same assertion holds over $\R$.
No separation of the old punctures modulo $p$ is required.

Now put $\mathcal C=\{0,1,2/3,-1/3\}$ and take any
rational $R\notin\mathcal C$.
For the pointed cover
\[
       C_R^{(8)}:\quad T-a=(R-a)X_a^8
                                      \quad(a\in\mathcal C),
\tag{A.125.2}
\]
every common octic character for a rational new puncture
$c\notin\mathcal C\cup\{R\}$ has zero normalized
evaluation at every place:
\[
                 \left(\frac{T-c}{R-c},\chi\right)=0.
\tag{A.125.3}
\]
Consequently, for every rational $c\notin\mathcal C$,
every finite set of places and one prescribed ordinary
point of (A.125.2) at each, avoiding $T=c$, a rational
coefficient $a_c\ne0$ exists such that
\[
 \begin{gathered}
 T-a=(R-a)X_a^8\quad(a\in\mathcal C),\qquad
 T-c=a_cY^8,\\
        (T,(X_a)_a,Y)\longmapsto(T,(X_a)_a)
 \end{gathered}
\tag{A.125.4}
\]
has an ordinary point over every completion and lifts
every prescribed old point.
This is finite exponent-eight initialization on the
actual degree-three Belyi source of A.114.
\end{statement}
\begin{proof}
\emph{The exact two-point equations for arbitrary old coordinates.}
Let $P_R=(R,\xi_i)$ and $P_T=(T,\eta_i)$ be two old
points over an odd completion $k=\Q_p$, with $R,T\ne c$.
For $T=R$ their evaluations are equal.
Otherwise define
\[
 x=\frac{T-c}{R-c},\qquad w_i=\eta_i/\xi_i,
       \qquad d_i=D_i/\xi_i^8.
\tag{A.125.5}
\]
Then
\[
             w_i^8=x+(1-x)d_i,
             \qquad d_i=\frac{w_i^8-x}{1-x}.
\tag{A.125.6}
\]
The individual quadratic characters have the exact
norm identity (A.118.13), with $w_i^2$ there replaced
by $w_i^4$.
The map of old covers $X_i\mapsto X_i^2$ and A.118
show that every character of $B_4$ has zero normalized
evaluation at these two points.

The first two cases of A.121, in which $B_8=B_4$, are
therefore finished.
In its remaining case one has
\[
 \mu_8\subset B_4,\qquad i\in K:=\Q(\sqrt{D_i}),
                    \qquad (x,-1)=(x,2)=0.
\tag{A.125.7}
\]
The inclusion $i\in K$ follows also from A.96: without
it $B_4$ would equal $K$.
Every additional character in A.121 comes from a
mixed relation, with $t_i\in\{0,1,2,3\}$,
\[
 \begin{gathered}
 \prod_iD_i^{t_i}=\sigma r_t^4,\qquad
 \prod_i d_i^{t_i}=\sigma r^4,\\
 r=r_t/\prod_i\xi_i^{2t_i},\qquad \sigma\in\{1,-1\}.
 \end{gathered}
\tag{A.125.8}
\]
The factor relating $r$ to $r_t$ is a square in $k$.
Thus its actual restricted character is $\chi_r$ for
$\sigma=1$, and $\chi_{16}+\chi_r$ for $\sigma=-1$.
Here $\chi_r$ is the quadratic character of $k(\sqrt r)$.
This keeps every positive and negative relation vector
in the complete list A.121.

If necessary make the exact replacement
\[
 x'=x^{-1},\qquad w_i'=w_i^{-1},\qquad
 d_i'=d_i/w_i^8,\qquad r'=r/\prod_iw_i^{2t_i}.
\tag{A.125.9}
\]
Equations (A.125.6)--(A.125.8) still hold.
The quadratic character of $r$ is unchanged, the
first entry is inverted, and the local squareclass
span of the $d_i$ is unchanged.
Consequently vanishing is equivalent for the new data.
We may take
$x=p^h u$ with $h\ge0$ and $u$ a unit, and write
$s_i=8v_p(w_i)\in8\Z$.
For $h>0$, the denominator $1-x$ in (A.125.6) is
a unit with residue one.
For $h=0$ and nonsquare $x$, its residue is again
nonzero, since the nonsquare $\bar u$ cannot equal one.

The cyclotomic character in A.121 is unramified at
every odd prime.  With the arithmetic Frobenius
convention, its actual values are
\[
\begin{array}{c|rrrrrrrr}
 p\bmod16&1&3&5&7&9&11&13&15\\ \hline
 4\chi_{16}(p)&0&3&1&2&2&1&3&0.
\end{array}
\tag{A.125.10}
\]
This follows directly from $\chi_{16}(5)=1/4$ and
$\chi_{16}(-1)=0$.
Its invariant on $x$ is $h\chi_{16}(p)$.
We use the actual local reciprocity map and its norm
kernel as in \cite[Theorem 4.1.2]{BIR_KedlayaCFT}.

\emph{Odd primes congruent to one modulo eight.}
Here $\mu_8\subset k$.
When $x$ is a nonsquare, $L=k(\alpha)$ with
$\alpha^8=x$ is cyclic of degree eight: it is its own
splitting field, and any proper subgroup of its group
in $\mu_8$ would fix $\alpha^4$, making $x$ a square.
The literal norms are
\[
 d_i=N_{L/k}\left(\frac{w_i-\alpha}{1-\alpha}\right).
\tag{A.125.11}
\]
For a positive relation $r^4$ is therefore a norm.
Its reciprocity equation puts the image of $r$ in the
subgroup $2\Z/8$, which is precisely the kernel of
restriction to the unique quadratic subfield
$k(\sqrt x)$.
Thus $(x,r)=0$.

For a negative relation fix a primitive $\zeta_8\in k$.
Now $(\zeta_8r)^4=-r^4$ is a norm, so the same argument
gives $(x,\zeta_8r)=0$.
The actual restriction of $\chi_{16}$ to $G_k$ is the
quadratic character of $k(\sqrt{\zeta_8})=k(\mu_{16})$,
possibly the zero character: this is immediate from
the $p=1,9$ entries of (A.125.10).
Hence the two terms of
$(x,\chi_{16}+\chi_r)$ cancel.
If $x$ is a square, its quadratic evaluations are zero
and $h$ is even; (A.125.10) makes its cyclotomic
evaluation zero as well.

\emph{Odd primes congruent to five modulo eight.}
The extension $k(\sqrt2)/k$ is unramified quadratic:
Frobenius sends
$\zeta_8+\zeta_8^{-1}=\sqrt2$ to its negative.
The second norm in (A.125.7) forces $h$ to be even.
Suppose first that $x$ is a nonsquare, so $\bar u$ is
a nonsquare and $k(\sqrt x)/k$ is unramified quadratic.

When $h\in8\Z$, the denominator in (A.125.6) is a
unit.  At $h=0$ this uses $\bar u\ne1$.
Even when $s_i=h$, the leading unit of $w_i^8$ cannot
cancel the nonsquare leading unit of $x$.
It follows that every $v_p(d_i)=\min(s_i,h)$ is a
multiple of eight.
Equation (A.125.8) gives $v_p(r)$ even, and both the
quadratic and cyclotomic evaluations are zero.

For $h\notin8\Z$, one has $h>0$ and $s_i\ne h$.
Let $\eta:\F_p^\times\to\mu_4$ be a quartic residue
character, chosen using any generator of the cyclic
residue-field multiplicative group.
Then $\eta(-1)=-1$ and $\eta(\bar u)=\pm i$.
For $s_i<h$, the normalized unit of $d_i$ is an
eighth-power residue; for $s_i>h$, it is $-\bar u$.
Put $b=\sum_{s_i>h}t_i$.
The normalized-unit and valuation parts of
(A.125.8) give the exact implications
\[
\begin{array}{c|c|c}
 \sigma&b\pmod4&v_p(r)\pmod2\\ \hline
 1&0&0\\
 -1&2&h/2,
\end{array}
 \qquad 4v_p(r)\equiv hb\pmod8.
\tag{A.125.12}
\]
Indeed the unit equation is
$\eta(-\bar u)^b=1$ or $-1$, respectively, and
$\eta(-\bar u)$ has order four.
For a positive relation (A.125.12) makes $r$ a norm
from the unramified quadratic extension.
For a negative relation its quadratic invariant is
$v_p(r)/2=h/4$ modulo $\Z$.
The cyclotomic invariant is also $h/4$ modulo $\Z$,
because $\chi_{16}(p)=1/4$ or $3/4$ and $h$ is even.
Their sum is zero.

If $x$ is a square, the positive characters already
vanish.  The negative character also vanishes whenever
$4\mid h$.
For the remaining possibility $h\equiv2\pmod4$, the
same quartic residue equation has
$\eta(-\bar u)=\pm1$ and right side $-1$.
It therefore forces $b$ odd.
But the valuation equation
$4v_p(r)=\sum_{s_i<h}t_is_i+hb$ would then be $2$
modulo four, a contradiction.
Thus no negative mixed relation occurs in that last
configuration.

\emph{Odd primes congruent to three modulo four.}
Now $k(i)/k$ is unramified quadratic, so the first
norm in (A.125.7) again forces $h$ even.
For nonsquare $x$, $\bar u$ is a nonsquare.
If $h\notin8\Z$, then $h>0$ and no $s_i$ equals $h$.
Each $d_i$ has even valuation and square normalized
unit: the two possible units are an eighth power or
$-\bar u$, and $-\bar u$ is a square since both
$-1$ and $\bar u$ are nonsquares.
Thus every $d_i$ is a square in $k$.
This contradicts the actual inclusion $i\in K$ in
(A.125.7), because
\[
 kK=k(\sqrt{D_i}:i)=k(\sqrt{d_i}:i)=k
\tag{A.125.13}
\]
would otherwise contain $k(i)\ne k$.

For $h\in8\Z$, the same no-cancellation calculation
as above gives $v_p(d_i)\in8\Z$ and $v_p(r)$ even.
The quadratic invariant for nonsquare $x$ is then
zero, as is $h\chi_{16}(p)$.

For square $x$ the positive characters vanish.
A negative relation with $h\equiv2\pmod4$ is
impossible: using the quadratic residue character
in (A.125.8) gives $(-1)^b=-1$, hence $b$ odd,
again contradicting the valuation equation modulo
four.
For $4\mid h$, (A.125.10) gives zero cyclotomic
invariant.  This completes every odd-prime case,
including arbitrary bad-reduction configurations
of the old punctures.

Over $\R$, all positive old eighth-power ratios put
$R,T$ in the same interval relative to the old
punctures.
If $c$ is outside the interval between them, $x>0$
and every real cyclic evaluation is zero.
If $c$ lies between them, all $D_i>0$, so the actual
boundary fibre has a real root tuple.
Its local stabilizer is the whole real Galois group;
every common character is therefore locally zero.
This proves the first assertion of the theorem.

\emph{The exact common-character maps at the dyadic place.}
Return to the four punctures $\mathcal C$ in (A.125.2).
The coordinate-square map $X_a\mapsto X_a^2$ sends
each old octic point to the quartic cover in A.120.
Thus every character of $B_4$ has zero normalized
evaluation.
We again need only the remaining case
$\mu_8\subset B_4$, so $(f,-1)=(f,2)=0$.
Every character of $B_8$ has order dividing four by
A.121; it kills actual fourth powers.

The valuation partition (A.120.11) applies because
every eighth power is a fourth power.
It excludes every change between distinct boundary
rows using the two explicit residues (A.120.12)--(13).
If the old valuation coordinates agree, set
$M=\max_a v_2(R-a)$.
The eighth-power unit equations give
$v_2(T-R)\ge M+5$.
Since $i\in K$, a negative rational-square subset
of the $D_a$ exists.
As in the equal-coordinate argument of A.123,
this forces $v_2(c-R)\le M+1$; otherwise every
$D_a$ would be $1$ modulo four.
Hence $f\in1+16\Z_2$ is an actual fourth power,
and every common-character evaluation vanishes.

For different depths in one finite boundary row
near $a$, put
$\delta=\max_{b\ne a}v_2(a-b)$.
The same neighbour ratio used in (A.120.14) is now
an eighth-power unit, so
\[
       \min\{v_2(R-a),v_2(T-a)\}\ge\delta+5.
\tag{A.125.14}
\]
The actual affine map $z\mapsto(z-a)/2^\delta$
divides the old coefficients by $2^\delta$ and
keeps every root coordinate.
It leaves each boundary radicand and $f$ unchanged.
The distinguished boundary becomes zero, all other
finite boundaries $b$ have $v_2(b)\le0$, and both
old depths are at least five.

For the infinite row with different depths, the
ratio in (A.120.15) is an eighth-power unit, so
both $-v_2(R)$ and $-v_2(T)$ are at least five.
Use the literal eighth-root cover isomorphism
\[
 U=1/T,\qquad Y_0=X_0^{-1},\qquad
 Y_{1/b}=X_b/X_0\quad(b=1,2/3,-1/3).
\tag{A.125.15}
\]
Its inverse is the one in (A.120.16).
On the rational boundary-fibre algebras the same
formulas give
\[
 D'_0=D_0^{-1},\qquad D'_{1/b}=D_b/D_0,\qquad
 f'=fR/T=fX_0^{-8}.
\tag{A.125.16}
\]
Consequently the actual common field $B_8$ and each
of its characters are identified by this algebra
isomorphism.  The same map commutes with coordinate
squaring, so it also identifies $B_4$ and its norm
conditions.
All mixed characters are transported as these
literal characters of the same field; none is
discarded by a change of its relation vector.
Since $X_0^{-8}$ is a norm from every cyclic field
of degree dividing four, the evaluation is unchanged.
The new finite punctures are $0,1,3/2,-3$, and
the two old points have depths at least five at zero.

\emph{The general deep-disc octic calculation.}
We have reached an old set consisting of zero and
nonzero $b$ with $v_2(b)\le0$, with
\[
 e=v_2(R),\quad e_T=v_2(T),\quad
 e,e_T\ge5,\quad e\ne e_T,\quad e\equiv e_T\pmod8.
\tag{A.125.17}
\]
The common field still satisfies $\mu_8\subset B_4$,
and all $B_4$ evaluations are still zero.
Write $s=\min(e,e_T)=v_2(T-R)$ and $k=v_2(c)$.
For $k\le0$, $f\in1+32\Z_2$ and is a fourth power.
For $1\le k\le4$, put $d=s-k$.
If $d\ge4$, the same conclusion holds.
The two known norm conditions $(f,-1)=(f,2)=0$
exclude $d=1,2$: a unit norm from $\Q_2(i)$ is
$1$ modulo four, and a unit norm from
$\Q_2(\sqrt2)$ is $1$ or $7$ modulo eight, by the
explicit quadratic forms in A.123.
Their intersection consists of units $1$ modulo
eight, whereas $d=1,2$ gives residues $3$ modulo
four or $5$ modulo eight, respectively.

For $d=3$, $f$ is a square and $k=s-3\ge2$.
In this case
\[
       v_2(D_0)=k-e\equiv-3\pmod8,\qquad
       D_b=\frac{1-c/b}{1-R/b}\in1+4\Z_2.
\tag{A.125.18}
\]
Any relation vector $t$ in A.121 has product
valuation divisible by four.
The odd first valuation in (A.125.18) therefore
forces $t_0=0$ in its specified range
$\{0,1,2,3\}$.
The remaining product is in $1+4\Z_2$ and cannot
be a negative fourth power.
Thus every relation in the complete A.121 list is
positive in this configuration.
That exact field formula makes $B_8$ multiquadratic,
and the square $f$ kills every character.

Finally take $k\ge5$.
Every $D_b$, $b\ne0$, belongs to $1+32\Z_2$ and
has an actual eighth root $u_b\in\Q_2$.
For any common character, its actual local cyclic
field $K_\chi$ lies in every field factor of the
boundary algebra: each local root stabilizer is
contained in the corresponding global stabilizer
on which the character vanishes.
Specializing the other roots gives the explicit
surjection
\[
 \mathcal E_{8,c}\otimes\Q_2
       \longrightarrow\Q_2[Z_0]/(Z_0^8-D_0),
                         \qquad Z_b\longmapsto u_b.
\tag{A.125.19}
\]
Every factor $E$ of its target still contains
$K_\chi$.
Using one root $\theta_E$ in each factor, the exact
polynomial norm identity reads
\[
 f=\frac{X_0^8-D_0}{1-D_0}
   =\prod_E N_{E/\Q_2}
                    \left(\frac{X_0-\theta_E}{1-\theta_E}\right).
\tag{A.125.20}
\]
Each factor is a norm from $K_\chi$ by transitivity.
The numerator and denominator are nonzero because
$T,c,R$ are distinct where required.
This proves every remaining dyadic evaluation.

The general odd and real assertions proved first
now apply to (A.125.2), including the odd prime
three where the five boundary points do not have
distinct reductions.
They complete (A.125.3).
Subtracting the global constant $(R-c,\chi)$ and
using reciprocity makes every sum in the exact
finite criterion A.85 zero.
That theorem gives the coefficient and actual
forgetful morphism (A.125.4), with the prescribed
finite collection of old points.
For $c=R$, the boundary has the rational root tuple
$(1,1,1,1)$ and common abelian field $\Q$; the
same criterion has no nonzero character to test.
This proves finite exponent-eight initialization.
No compatible choice through different exponents
or infinite collection of prescribed old points
is asserted.
\end{proof}

The exact norm, residue, mixed-relation and source
geometry checks are retained in
\path{code/cubic_source_eighth_certificate.py}
and its output.  These finite checks supplement
the unbounded local and finite-initialization proof.
